\documentclass[10pt]{article}
\usepackage[margin=1.5cm]{geometry}
\usepackage{amsmath}
\usepackage[spanish]{babel}
\usepackage{parskip}
\pagestyle{empty}

\begin{document}

% Evaluación 1
\begin{minipage}[t][0.23\textheight][t]{0.48\textwidth}
\textbf{Evaluación 1 – Mecánica de los Fluidos} \hfill Nombre: \hrulefill

\vspace{1mm}
\textbf{Curso: 5º Año – Técnica} \\
\textbf{Duración: 80 minutos}

\vspace{2mm}
\textbf{1.} ¿Cuál es la presión absoluta en el fondo de un tanque de agua de $6,5\,\text{m}$ de profundidad?

\vspace{2mm}
\textbf{2.} ¿Qué presión se transmite a todos los puntos del líquido si se aplica $F=25\,\text{N}$ sobre un émbolo de $A=5\,\text{cm}^2$?

\vspace{2mm}
\textbf{3.} Una prensa hidráulica tiene $A_1=0{,}002\,\text{m}^2$ y $A_2=0{,}1\,\text{m}^2$. Si se aplica $F_1=100\,\text{N}$, calcular $F_2$.

\vfill
\textit{Mostrar procedimiento.}
\end{minipage}
\hfill
% Evaluación 2
\begin{minipage}[t][0.23\textheight][t]{0.48\textwidth}
\textbf{Evaluación 2 – Mecánica de los Fluidos} \hfill Nombre: \hrulefill

\vspace{1mm}
\textbf{Curso: 5º Año – Técnica} \\
\textbf{Duración: 80 minutos}

\vspace{2mm}
\textbf{1.} Determinar la presión a $10\,\text{m}$ de profundidad en agua (densidad $1000\,\text{kg/m}^3$).

\vspace{2mm}
\textbf{2.} ¿Qué fuerza hay que aplicar sobre un émbolo de $2\,\text{cm}^2$ para generar una presión de $3\times10^5\,\text{Pa}$?

\vspace{2mm}
\textbf{3.} En una prensa hidráulica, el plato mayor tiene $r=0{,}25\,\text{m}$ y el menor $r=0{,}05\,\text{m}$. Calcular $F_1$ para levantar $P=1500\,\text{N}$.

\vfill
\textit{Mostrar unidades y pasos.}
\end{minipage}

\vspace{4mm}

% Evaluación 3
\begin{minipage}[t][0.23\textheight][t]{0.48\textwidth}
\textbf{Evaluación 3 – Mecánica de los Fluidos} \hfill Nombre: \hrulefill

\vspace{1mm}
\textbf{Curso: 5º Año – Técnica} \\
\textbf{Duración: 80 minutos}

\vspace{2mm}
\textbf{1.} ¿Cuál es la diferencia de presión entre dos puntos separados $4\,\text{m}$ en agua?

\vspace{2mm}
\textbf{2.} En un sistema cerrado se aplica $F=50\,\text{N}$ sobre $A=2{,}5\,\text{cm}^2$. ¿Cuál es la presión en Pa?

\vspace{2mm}
\textbf{3.} Si el radio del plato mayor es el triple del menor ($r_1 = 3 r_2$), y se aplica $F_2=5\,\text{N}$, calcular $F_1$.

\vfill
\textit{Justificar resultados con fórmulas.}
\end{minipage}
\hfill
% Evaluación 4
\begin{minipage}[t][0.23\textheight][t]{0.48\textwidth}
\textbf{Evaluación 4 – Mecánica de los Fluidos} \hfill Nombre: \hrulefill

\vspace{1mm}
\textbf{Curso: 5º Año – Técnica} \\
\textbf{Duración: 80 minutos}

\vspace{2mm}
\textbf{1.} ¿Qué altura de agua genera una presión manométrica de $1{,}5\times10^5\,\text{Pa}$?

\vspace{2mm}
\textbf{2.} Aplicando $F=20\,\text{N}$ sobre un émbolo de $4\,\text{cm}^2$, ¿cuál es la presión generada?

\vspace{2mm}
\textbf{3.} Se coloca una masa de $m=6\,\text{kg}$ en el émbolo menor. ¿Qué masa se puede levantar si el émbolo mayor tiene el triple de radio?

\vfill
\textit{Responder con claridad y orden.}
\end{minipage}

\end{document}
